\section{Sixth Wound: The Saturated Shadow}

The five wounds above were handed to me. Most are the historical objections, raised by people who wanted the mechanical program to succeed, and every one of them is older than I am. This one nobody handed me. I found it while writing the section that closes this chapter, and I have put it here, in the ledger with the others, rather than leaving it where I found it.

Begin with something the shadow push has needed from the start and has so far been allowed to have quietly. Gravity, in this sketch, is a push from a rain that is blocked: a body pushes its neighbour because it removes some of the flux that would otherwise have arrived from that direction. For the push to come out proportional to \emph{mass} --- which is what we observe, and to great precision --- each massion must have only a very small chance of being stopped by any given piece of matter. Matter has to be nearly transparent to the crowd. If it were not --- if a body stopped a large fraction of what reached it --- the far side of the body would sit in the near side's shadow and contribute nothing, and the push would stop tracking the mass and begin tracking the \emph{silhouette}.

Ordinary matter is obliging about this. The Earth is a ball of rock eight thousand miles thick and its gravity still counts every atom, which is a direct measurement, if you like, of how nearly transparent matter is to whatever gravity is made of. Fine. But now ask about a body that is not obliging: something collapsed, something dense enough to trap light. Is it still transparent to the crowd?

If it is not, the arithmetic is brutal. A collapsed body's radius grows in proportion to its mass, so its silhouette --- the area it blocks --- grows as the \emph{square} of its mass. An opaque body would therefore push as the square of its mass, while everything we can weigh insists the push grows as the mass itself. The gap is not a rounding error. Between a hole of ten suns and the one at the centre of this galaxy the two rules differ by a factor of hundreds of thousands. And nature has been unkind about where it set the examination: collapsed bodies are precisely where we now measure gravity best --- stars whipping around the dark mass at our galaxy's centre, the ringed shadow the Event Horizon Telescope photographed, the waveforms of inspiralling pairs matching general relativity to many decimal places. If the shadow saturates there, that agreement is a coincidence, and I do not believe in coincidences of that size.

The best reply I have is real but it is not free. A horizon is our layer's \emph{optical} fact: it is defined by what our light cannot climb out of. The crowd is not our light. It belongs to the layer below, it moves very much faster, and it is under no obligation to be stopped by whatever stops a photon. A collapsed body may therefore be black and yet nearly transparent --- opaque to us, and hardly there at all as far as the rain is concerned. Then the push tracks the mass, the orbits come out right, and the wound closes.

But it closes by paying elsewhere, and I want the reader to see the coin leave my hand. The next-to-last section of this chapter will note that black holes carry an entropy proportional to the \emph{area} of the horizon, and will offer the tower an explanation of it. There are two explanations available. The cheap one asks only that our layer be unable to see past the horizon, and it survives this reply intact. The satisfying one --- that the crowd itself is stopped at the surface, so that the shell is genuinely all there is to interact with, the interior having been pressed dense by the shadow of its own shell --- is exactly the opacity this reply has just denied. So the dial has two ends and neither is comfortable. Turn it to opaque: the area law becomes mechanical and gravity goes wrong by a factor of hundreds of thousands. Turn it to transparent: gravity is saved and the area law falls back to bookkeeping. It is the first wound's shape all over again --- a pincer, with the same two jaws closing on the same one knob --- and it is the newest and least settled thing in this book.

What would settle it is a number: the massion's chance of being stopped, per unit of matter it passes through, at densities far beyond anything we can make. Everything above turns on where that number puts the saturation, and I do not know where it puts it. That is the sixth problem of the research program below, and of the six it is the one I would most like to see somebody take away from me.

